\documentclass[10pt,a4paper]{amsart}
\usepackage[margin=19mm]{geometry}
\usepackage{amsmath,amssymb,mathtools}
\usepackage[scr=rsfs,cal=euler]{mathalfa}
\usepackage{xfrac}
\usepackage[unicode,hidelinks,bookmarks=false]{hyperref}
\newcommand{\ie}{i.e.\ }
\newcommand{\cf}{cf.\ }
\newcommand{\resp}{respectively\ }
\newcommand{\Subsection}{\subsection}
\providecommand{\nobreakdash}{\nobreak\hbox{-}}
\setlength{\emergencystretch}{2em}
\multlinegap0pt
\usepackage{amssymb}
\newtheorem{lemma}{Lemma}
\title{On the convergence of a perturbed one dimensional Mann's process}
\author{Ramzi May}
\date{Source: arXiv:2504.16154v2 (CC BY 4.0)}
\begin{document}
\maketitle

\noindent\emph{Source and licence.} On the convergence of a perturbed one dimensional Mann's process. Version arXiv:2504.16154v2 (CC BY 4.0), distributed by arXiv under the Creative Commons Attribution 4.0 International licence, http://creativecommons.org/licenses/by/4.0/, as recorded in the arXiv licence metadata for this exact version and shown on the abstract page https://arxiv.org/abs/2504.16154v2. Full licence notice: \texttt{licenses/arxiv-2504.16154-CC-BY-4.0.txt}.\par\smallskip

\noindent\textit{Editorial scope.} Counted unit: the article's lemma l2 (source0.tex lines 244--256). The article's complete original proof of it is delivered with it (source0.tex lines 296--360, one \texttt{proof} environment of 65 lines, which the article places away from the statement). No mathematics is written, repaired or re-derived: the statement and the proof are the article's own lines, in the article's own notation, and no retained line is substituted or re-flowed.  Retained uncounted as the source-local context the proof needs: lines 260--273.  Nothing was omitted from any retained span, so the package carries no omission record.  No retained line contains a citation, so the wrapper carries no bibliography and no bibliography entry.  No retained line contains a figure environment, an image-inclusion command or drawing code, so no image is delivered and the package declares no asset dependency.  Editorial formatting only, in addition: an AMS \texttt{amsart} preamble that restates the article's own declarations and macros used by the retained lines, namely the environments \texttt{lemma} on separate counters, together with the standard packages those lines need.  The retained environments print in this document as Lemma 1 in order of appearance, whereas the article prints the same items with the numbering of its own full document.  The complete native source of the article as distributed in the arXiv e-print for this exact version is bundled unchanged as \texttt{sources/176920/source0.tex}.
\begin{lemma}
\label{l3}
let $u:[0,\infty )\to \mathbb{R}$ be a continuous and bounded
function. If $\omega (u(t))=[a,b]$ with $a<b$ then for any real numbers $c$
and $d$ such that $a<c<d<b$ and any $T_{0}>0$ there exist four reals numbers
$s_{1},s_{2},\tau _{1}$ and $\tau _{2}$ such that:
\begin{enumerate}
\item $\tau _{i}>s_{i}>T_{0}$, $i=1,2.$
\item $u(t)\in [c,d]$ for any $t\in [s_{i},\tau _{i}],$ $%
i=1,2.$
\item $u(s_{1})=c$ and $u(\tau _{1})=d.$
\item $u(s_{2})=d$ and $u(\tau _{2})=c.$
\end{enumerate}
\end{lemma}

\begin{lemma}
\label{l2}
let $u:[0,\infty )\to \mathbb{R}$ be a continuous and bounded
function. Then
\begin{enumerate}
\item $\omega (u(t))=\cap _{k\in \mathbb{N}}\overline{u([k,\infty ))}^{\mathbb{R}%
},$ where $\overline{u([k,\infty ))}^{\mathbb{R}}$ denotes the closure of
the set $u([k,\infty ))$ in $\mathbb{R}.$
\item $\omega (u(t))$ is a nonempty compact interval of $\mathbb{R}.$
\item $u(t)$ converges to some real number $u_{\infty }$ as $%
t\rightarrow \infty $ if ad only if $\omega (u(t))=\{u_{\infty }\}.$
\end{enumerate}
\end{lemma}

\begin{proof}
From Lemma \ref{l2}, there exist two real numbers $0\leq a\leq b\leq 1$ such that $\omega
(x(t))=[a,b].$ First, for the sake of a contradiction, we suppose that $a<b$.  Now let $w$ be an arbitrary element of the open interval $(a,b).$ Let us
suppose that $f(w)>w.$ Then there exist $\varepsilon ,\delta >0$ such that $%
[w-\delta ,w+\delta ]\subset (a,b)$ and $f(z)>z+\varepsilon $ for any $z\in [w-\delta ,w+\delta ].$ Let $T_{0}>0$ such that $r(t)\geq
-\varepsilon \theta (t),\forall t\geq T_{0}.$ According to Lemma \ref{l3}, there
exist $\tau _{2}>s_{2}>T_{0}$ such that $x(s_{2})=w+\delta ,~x(\tau
_{2})=w-\delta ,$ and $x(t)\in [w-\delta ,w+\delta ]$ for any $t\in
[s_{2},\tau _{2}].$ Therefore, for every $t\in [s_{2},\tau
_{2}],$%
\begin{eqnarray*}
x^{\prime }(t) &=&\theta (t)(f(x(t))-x(t))+r(t) \\
&\geq &\varepsilon \theta (t)+r(t) \\
&\geq &0.
\end{eqnarray*}%
Hence by integrating on the interval $[s_{2},\tau _{2}],$ we get the
contradiction $-2\delta \geq 0.$ Similarly, by using the times $s_{1}$ and $%
\tau _{1}$ instead of $s_{2}$ and $\tau _{2},$\ we can verify that the
assumption $f(w)<w$ leads also to a contradiction. We then conclude that if $%
a<b$ then $f(w)=w$ for any $w\in [a,b].$ We continuous working under
the assumption $a<b.$ Let $c=\frac{a+b}{2}$ and $\alpha =\frac{b-a}{4}.$ Let
$T_{1}>0$ be greater enough such that
\begin{equation}
\left\vert \int_{s}^{\tau }r(t)dt\right\vert \leq \alpha ,~\forall \tau
>s>T_{1}.  \label{h5}
\end{equation}%
According to Lemma \ref{l3}, there exist $\tau _{1}>s_{1}>T_{1}$ such that $%
x(s_{1})=c-\alpha ,~x(\tau _{1})=c+\alpha ,$ and $x(t)\in \lbrack c-\alpha
,c+\alpha ]$ for any $t\in \lbrack s_{1},\tau _{1}].$ Therefore, for any $%
t\in \lbrack s_{1},\tau _{1}],$

\begin{eqnarray*}
x^{\prime }(t) &=&\theta (t)(f(x(t))-x(t))+r(t) \\
&=&r(t).
\end{eqnarray*}%
Hence, by integrating the last equality between $s_{1}$ and $\tau _{1}$ we
get
\[
2\alpha =\int_{s_{1}}^{\tau _{1}}r(t)dt,
\]%
which contradicts (\ref{h5}). We therefore conclude that $a=b$ which means,
thanks to Lemma \ref{l2}, that $x(t)$ converges as $t\rightarrow \infty $ to some the
real number $x_{\infty }=a=b.$ It is clear that $x_{\infty }\in
\lbrack 0,1],$ let us show that it is a fixed point of $f.$ For the sake of
absurdity, we assume for instance that $f(x_{\infty })<x_{\infty }.$ The
continuity of $x(t)$ and $f$ and the assumption on $r(t)$ ensure the
existence of $\varepsilon >0$ and $T_{2}$ such that for any $t\geq T_{2}$%
\begin{eqnarray*}
x(t)-f(x(t)) &\geq &2\varepsilon , \\
r(t) &\leq &\varepsilon \theta (t).
\end{eqnarray*}%
Therefore, for every $t\geq T_{2},$
\begin{eqnarray*}
x^{\prime }(t) &=&\theta (t)(f(x(t))-x(t))+r(t) \\
&\leq &-\varepsilon \theta (t).
\end{eqnarray*}%
Integrating the last inequality between $T_{2}$ and $T>T_{2}$ and then letting $%
T\rightarrow \infty ,$ we get the inequality%
\[
\varepsilon \int_{T_{2}}^{\infty }\theta (t)dt\leq x(T_{2})-x_{\infty }
\]%
that contradict the assumption on $\theta .$ Similarly, we can show that
$f(x_{\infty })$ can not be greater than $x_{\infty }.$ We therefore
conclude that $f(x_{\infty })=x_{\infty }.$
\end{proof}

\end{document}
